\addcontentsline{toc}{subsubsection}{Solution 45. The Conjugate-Radical Lock}

\begin{dlsummary}[title={\bfseries Solution 45. The Conjugate-Radical Lock},breakable]

\begin{proof}
The unique solution is
\[
\boxed{
f(x)=\frac{(2+\sqrt3)^x+(2-\sqrt3)^x}{2}
\qquad(x\ge0).
}
\]

\paragraph{1. Remove the radicals.}

For \(t\ge1\), define
\[
H(t)=t+\sqrt{t^2-1}.
\]
This function is strictly increasing, and
\[
H(t)^{-1}=t-\sqrt{t^2-1}.
\]

For \(a,b\ge1\), direct expansion gives
\[
H\!\left(ab+\sqrt{a^2-1}\sqrt{b^2-1}\right)=H(a)H(b).
\tag{3}
\]
Also,
\[
H(2a^2-1)=H(a)^2. \tag{4}
\]

Set
\[
g(x)=\log H(f(x)).
\]
Then \(g(x)\ge0\). Applying \(H\) and taking logarithms in (1) and (2), using (3) and (4), gives
\[
g(x+y)\ge g(x)+g(y),
\qquad
g(2x)\le2g(x).
\]
The first inequality with \(y=x\) supplies the reverse doubling bound. Therefore
\[
\boxed{g(2x)=2g(x).} \tag{5}
\]

\paragraph{2. Recover monotonicity and rational values.}

Since \(g\) is nonnegative and superadditive, it is nondecreasing:
\[
b>a\ge0
\implies
g(b)\ge g(a)+g(b-a)\ge g(a).
\]

Write
\[
L=g(1)=\log(2+\sqrt3).
\]
Superadditivity gives \(g(n)\ge nL\) for every positive integer \(n\).

Choose a power of two \(N>n\). By (5), \(g(N)=NL\), so
\[
NL=g(N)\ge g(n)+g(N-n)\ge g(n)+(N-n)L.
\]
Thus \(g(n)\le nL\), proving
\[
g(n)=nL.
\]
Repeated doubling now yields
\[
g\!\left(\frac{m}{2^k}\right)=\frac{m}{2^k}L
\]
for all nonnegative integers \(m,k\).

\paragraph{3. Squeeze arbitrary real inputs.}

Every nonnegative real number can be approached from below and above by numbers of the form \(m/2^k\). Monotonicity therefore forces
\[
g(x)=Lx\qquad(x\ge0).
\]

Hence \(H(f(x))=(2+\sqrt3)^x\). Using the reciprocal identity for \(H\),
\[
2f(x)=H(f(x))+H(f(x))^{-1},
\]
which gives the claimed formula.

\paragraph{4. Verify.}

Put \(A=2+\sqrt3\), so \(A^{-1}=2-\sqrt3\). For the proposed function,
\[
\sqrt{f(x)^2-1}=\frac{A^x-A^{-x}}2
\qquad(x\ge0).
\]
Expansion shows that (1) holds with equality, as does (2). The domain, range, and normalisation are also satisfied.
\end{proof}
\end{dlsummary}

\begin{takeaways}
\begin{itemize}
\item \textbf{Use conjugate radicals:} The expressions \(t+\sqrt{t^2-1}\) and \(t-\sqrt{t^2-1}\) are reciprocals.
\item \textbf{Linearise before estimating:} A transformation can turn a radical addition rule into superadditivity.
\item \textbf{Exact doubling is powerful:} Together with nonnegativity and superadditivity, it determines values first on integers, then on dyadic rationals.
\item \textbf{Derive regularity:} Monotonicity justifies the final squeeze; continuity need not be assumed.
\end{itemize}
\end{takeaways}
